\section{Methodology} \label{sec:methodology}

To make automated MOSA integration amenable to formal verification and reinforcement learning while maintaining a systems-theoretic approach, we formulate interface mediation as an \emph{Integration Game}.  The framework defines a structured design space of positions (\emph{Mediation Boards}), legal moves (\emph{Mediation Operators}), winning conditions (typing, kinematic, and behavioral satisfiability), and quantitative trade-off objectives (\emph{Integration Cost Functionals}).  Every interface element is typed by a single Stateful \emph{Color}, characterized by the UDDL data semantic layers of Section~\ref{sec:bg-uddl} along with runtime representation, flow kinematics, and provenance; behavioral obligations are expressed as uLSC contracts equivalent to Wymorian $IOR$s; and search is guided by a Reward Machine compiled from those contracts.

\subsection{Elemental Decomposition} \label{sec:meth-isomorphism}

We decompose an integration problem as follows:

\begin{enumerate}[topsep=2pt, itemsep=1pt, leftmargin=1.8em]
    \item \textbf{Producers and Consumers.} Abstractions of discrete component systems whose interfaces are to be integrated.  Producer and Consumer are roles with respect to a particular exchange.  A component may produce in one exchange and consume in another.
    \item \textbf{Offers and Needs.} A Producer exposes \emph{Offers} $O_i \in \mathcal{O}$ and a Consumer declares \emph{Needs} $N_j \in \mathcal{N}$.
    \item \textbf{Messages.} Each Offer or Need is typed as a structured \emph{Message} which aggregates Fields as a record, a sequence, or a choice (Union).
    \item \textbf{Fields and Colors.} The atomic unit of mediated information is a \emph{Field}.  Each Field carries a Color $\kappa$ (defined below) and, at runtime, a value $v_{\text{val}}$.
    \item \textbf{Mediation Board.} A \emph{Mediation Board} $\mathcal{B} = (\mathcal{O}, \mathcal{N}, \mathcal{V}_{\text{op}}, \mathcal{E}_{\text{route}})$ is a directed mediation graph whose exterior boundary comprises Offers and Needs.  Its interior vertices $\mathcal{V}_{\text{op}}$ are \emph{Mediation Operators} drawn from the color-refined \emph{Technology Catalog} $TYR_{\mathcal{K}}$ (each operator itself a discrete Wymorian system $Z$) and may alter Color.  Its edges $\mathcal{E}_{\text{route}}$ are connections, each realized by a \emph{Route Carrier}: the transport medium that moves a Message between two ports (e.g., shared-memory segment, DDS topic, MIL-STD-1553 bus, etc.).  A Route Carrier represents transport cost and latency, denoted $c_{\text{carrier}}(r)$ for $r \in \mathcal{E}_{\text{route}}$.  Any wire encoding it requires is transparent to Messages/Fields and accounted for in $c_{\text{carrier}}$.
\end{enumerate}

Figure~\ref{fig:metamodel} summarizes the metamodel.

\begin{figure}[t]
\centering
\resizebox{\linewidth}{!}{%
\begin{tikzpicture}[
    >=Stealth,
    font=\footnotesize\sffamily,
    mbox/.style={draw=compBorder, fill=compFill, thick, rounded corners=3pt, minimum width=2.7cm, minimum height=1.0cm, align=center},
    offerbox/.style={draw=techNavy, fill=techNavy!8, thick, rounded corners=3pt, minimum width=2.7cm, minimum height=1.0cm, align=center},
    needbox/.style={draw=boundGreen, fill=boundGreen!10, thick, rounded corners=3pt, minimum width=2.7cm, minimum height=1.0cm, align=center},
    fieldbox/.style={draw=poiPurple, fill=poiPurple!8, thick, rounded corners=3pt, minimum width=3.6cm, minimum height=1.0cm, align=center},
    opbox/.style={draw=accentAmber!90!black, fill=accentAmber!10, thick, rounded corners=3pt, minimum width=2.7cm, minimum height=0.9cm, align=center},
    contains/.style={->, semithick, draw=black!55},
    carrier/.style={->, very thick, draw=fsrBlue},
    opedge/.style={->, thick, dashed, draw=accentAmber!90!black},
    elabel/.style={font=\scriptsize\sffamily, text=black!70}
]
    % Producer row
    \node[mbox] (za) at (0, 1.7) {\textbf{Producer}\\[0pt]{\scriptsize Component $Z_A$}};
    \node[offerbox] (offer) at (3.8, 1.7) {\textbf{Offer $O_i$}\\[0pt]{\scriptsize Producer role}};
    \node[mbox] (msg1) at (7.6, 1.7) {\textbf{Message}\\[0pt]{\scriptsize record / seq / choice}};
    \node[fieldbox] (fld1) at (11.8, 1.7) {\textbf{Field $f$}\\[0pt]{\scriptsize $\kappa = (C, L, P, \rho, k, \pi)$}};

    % Consumer row
    \node[mbox] (zb) at (0, -1.7) {\textbf{Consumer}\\[0pt]{\scriptsize Component $Z_B$}};
    \node[needbox] (need) at (3.8, -1.7) {\textbf{Need $N_j$}\\[0pt]{\scriptsize Consumer role}};
    \node[mbox] (msg2) at (7.6, -1.7) {\textbf{Message}\\[0pt]{\scriptsize record / seq / choice}};
    \node[fieldbox] (fld2) at (11.8, -1.7) {\textbf{Field $f'$}\\[0pt]{\scriptsize $\kappa' = (C', L', P', \rho', k', \pi')$}};

    % Mediation operator between Offer and Need
    \node[opbox] (op) at (3.8, 0) {\textbf{Operator $v \in \mathcal{V}_{\text{op}}$}\\[0pt]{\scriptsize may change $\kappa$}};

    % Containment edges
    \draw[contains] (za) -- node[above=1pt, elabel] {exposes} (offer);
    \draw[contains] (offer) -- node[above=1pt, elabel] {has} (msg1);
    \draw[contains] (msg1) -- node[above=1pt, elabel] {aggregates} (fld1);
    \draw[contains] (zb) -- node[above=1pt, elabel] {declares} (need);
    \draw[contains] (need) -- node[above=1pt, elabel] {has} (msg2);
    \draw[contains] (msg2) -- node[above=1pt, elabel] {aggregates} (fld2);

    % Route carriers (E_route): preserve kappa
    \draw[carrier] (offer.south) -- node[right=2pt, elabel, text=fsrBlue!90!black] {Route Carrier $e \in \mathcal{E}_{\text{route}}$} (op.north);
    \draw[carrier] (op.south) -- node[right=2pt, elabel, text=fsrBlue!90!black] {Route Carrier $e' \in \mathcal{E}_{\text{route}}$} (need.north);

    % Operator effect on Field Colors
    \draw[opedge] (fld1.south) -- node[right=3pt, elabel, text=accentAmber!70!black, align=left] {operators realize $\Delta\kappa$;\\carriers preserve $\kappa$} (fld2.north);
    \draw[opedge] (op.east) -- (11.8, 0);
\end{tikzpicture}%
}
\caption{Metamodel for the Integration Game: Component $\rightarrow$ Offer/Need $\rightarrow$ Message $\rightarrow$ Field typed by Color $\kappa$.  Solid blue edges are Route Carriers ($\mathcal{E}_{\text{route}}$), incurring transport cost $c_{\text{carrier}}$; dashed edges are Mediation Operators ($\mathcal{V}_{\text{op}}$), which may mutate $\kappa$ along one or more of its axes.}
\label{fig:metamodel}
\end{figure}

\paragraph{Color.} A Field is typed by its \emph{Color},
\begin{equation}
    \kappa = (C, L, P, \rho, k, \pi) \in \mathbb{C} \times \mathbb{L} \times \mathbb{P} \times \mathbb{R}_{\text{rep}} \times \mathbb{K} \times \Pi =: \mathcal{K},
    \label{eq:stateful-color}
\end{equation}
where $\mathcal{K}$ denotes the Color space.  Color is distinct from the runtime value $v_{\text{val}}$ it characterizes: a value ranges over a domain determined by $\kappa$, whereas $\kappa$ itself is fixed by the declaration of a port slot.  Color is stateful in three senses:
\begin{itemize}[topsep=2pt, itemsep=1pt, leftmargin=1.8em]
    \item \textbf{Static Coloring $(C, L, P, \rho)$.} $C \in \mathbb{C}$ identifies the Conceptual Entity and its rooted navigational path (e.g., \texttt{Aircraft.sensor.position}).  While the full UDDL standard specifies a rich SQL-like query grammar supporting multi-entity joins and projections, our implementation models $C$ via a simplified canonical semantic paths over CDM entity compositions and observables, which suffices to capture hierarchical $\Delta C$ while keeping conceptual matching decidable and lightweight.  $L \in \mathbb{L}$ is the Logical Frame (measurement system, frame, units), and $P \in \mathbb{P}$ the Platform realization (mirroring the UDDL hierarchy of Section~\ref{sec:bg-uddl}).  $\rho \in \mathbb{R}_{\text{rep}}$ is the in-memory, language-specific realization of the value (e.g., a C \texttt{double}, a CPython \texttt{PyObject*}, a JVM object, or a ROS\,2/DDS language binding).  Note that $\rho$ is not wire serialization, and must be contrasted with the Platform axis $P$, which is the UDDL PDM abstract machine type and precision (e.g., IEEE-754 \texttt{binary64}).  Two ports may align on $(C, L, P)$ (e.g., a WGS84 latitude as \texttt{binary64}) but still differ in $\rho$ (e.g., one has a \texttt{PyObject*} and the other a C-ABI \texttt{double}).  Thus, coupling them directly is unsafe, and a legal move must bind or marshal one realization into the other (Section~\ref{sec:meth-moves}).
    \item \textbf{Kinematics $k$.} $k \in \mathbb{K}$ records the spatiotemporal delivery profile at a boundary port: its topological location $x$ on the board, flow velocity $v$ (an active push flow with $v > 0$ versus a stationary pull potential with $v = 0$), and flow acceleration\footnote{A non-zero $a$ captures impulsive mode transitions, in which a store decelerates a push stream to rest ($a < 0$) or a polling pump accelerates a pull source into a periodic push stream ($a > 0$), while a rate differential $\Delta v$ across a multi-rate boundary dictates downstream buffering and staleness requirements.} $a = \Delta v / \Delta t$, written $k = (x, v, a)$, together with staleness (sample age and its admissible bound).
    \item \textbf{Provenance $\pi$.} $\pi \in \Pi$ records the source lineage and other properties such as assurance level of the Field, and is extended as operators transform it.  For example, a triple modular redundancy (TMR) voter fuses three redundant Fields and emits one whose $\pi$ records the vote history.
\end{itemize}

Luckily, port constraints expressed over $\mathcal{K}$ preserve the algebraic structure of Wymorian coupling.  For a port slot $p$, write $\kappa(p)$ for its Color and $A_{\kappa(p)}$ for the alphabet of values it admits; $A_{\kappa(p)}$ is determined by the projection $(C, L, P, \rho)$, while $k$ and $\pi$ constrain admissible delivery timing and lineage.  A \emph{Play} $\mathcal{P}$ is a candidate integration on the board; composition of its components induces the \emph{Coupled Resultant} $Z_b = \mathrm{RSY}(\mathcal{P})$.

\begin{proposition}[Color Typed Compatibility] \label{prop:wymore-iso}
Let $TYR$ be a catalog of discrete component systems and Mediation Operators whose ports are partitioned into elementary Field slots, and let $\kappa$ assign a Color $\kappa(p) \in \mathcal{K}$ to each slot $p$.  Define the color-refined Technology Catalog $TYR_{\mathcal{K}}$ by restricting the alphabet of each slot to $A_{\kappa(p)}$.  Then a Play $\mathcal{P}$ is legal (i.e., every wired port pair $(p_{\text{out}}, p_{\text{in}})$ satisfies $\mathrm{proj}_{CL P\rho}(\kappa(p_{\text{out}})) = \mathrm{proj}_{CL P\rho}(\kappa(p_{\text{in}}))$, the kinematics $k(p_{\text{out}})$ are admissible for $k(p_{\text{in}})$, and $\pi(p_{\text{out}})$ meets the provenance predicate of $p_{\text{in}}$) if and only if $\mathcal{P}$ forms a valid Wymorian System Coupling Recipe $SCR = (VSCR, CSCR)$ over $TYR_{\mathcal{K}}$ whose wired pairs additionally satisfy these kinematic and provenance admissibility conditions.  In that case $A_{\kappa(p_{\text{out}})} = A_{\kappa(p_{\text{in}})}$ for every pair in $CSCR$, and synchronous composition induces a unique Coupled Resultant discrete system $Z_b = \mathrm{RSY}(SCR) \in \mathrm{CTL}(TYR_{\mathcal{K}})$.
\end{proposition}

\begin{proof}[Proof Sketch]
In T3SD, an $SCR = (VSCR, CSCR)$ is well-formed iff $CSCR$ is a bijection between a subset of component output ports and a subset of component input ports such that the output alphabet of each source port equals the input alphabet of its destination~\cite{Wymore1993}.  Because $A_{\kappa(p)}$ is a function of the projection $(C, L, P, \rho)$ of $\kappa(p)$, equality of these projections implies $A_{\kappa(p_{\text{out}})} = A_{\kappa(p_{\text{in}})}$ and, up to declared equivalence of Colors, conversely.  In particular, two ports that agree on $(C, L, P)$ but differ in $\rho$ have distinct alphabets, so their direct coupling is not a well-formed $CSCR$ edge; this is the formal counterpart of the unsafe \texttt{PyObject*}/C-ABI couple above.  The axes $k$ and $\pi$ do not change alphabets but gate coupling: they restrict wiring to pairs whose delivery profile and lineage satisfy the Need's declared predicates, and are discharged as the Tier~1 and Tier~2 conditions of Section~\ref{sec:meth-winning}.  Route Carriers in $\mathcal{E}_{\text{route}}$ connect ports of identical Color and therefore are edges of $CSCR$ rather than members of $VSCR$.  Any Color discrepancy $\Delta\kappa$ between an Offer and a Need must be resolved by an interior Mediation Operator $v \in \mathcal{V}_{\text{op}} \subseteq VSCR$ whose input port has alphabet $A_{\kappa_{\text{in}}}$ and whose output port has alphabet $A_{\kappa_{\text{out}}}$.  The step transition $\delta_b$ and readout $\rho_{\text{out},b}$ of $Z_b = \mathrm{RSY}(SCR)$ then follow from the standard Wymorian coupling algebra.
\end{proof}

\subsection{Mediation Operators and Legal Moves} \label{sec:meth-moves}

Mediation Operators in $TYR_{\mathcal{K}}$ are organized into four families (Table~\ref{tab:legal-moves}):

\begin{enumerate}[topsep=2pt, itemsep=1pt, leftmargin=1.8em]
    \item \textbf{Stateful Recoloring:} Operators that transform a Field along one or more axes of $\kappa$.  A \emph{Conceptual Recolor} ($\Delta C$) maps a Field between equivalent or related CDM concepts; a \emph{Logical Recolor} ($\Delta L$) converts units, reference frames, or coordinate systems via a declared quantitative transformation; a \emph{Platform Recolor} ($\Delta P$) widens, narrows, or repacks data types across PDM encodings; and a \emph{Runtime Representation Binding/Marshaling} operator ($\Delta\rho$) re-expresses the same value in a different in-memory or language-runtime realization (foreign function interfaces, JNI bridges, language-binding conversions, etc.).
    \item \textbf{Structural Aggregation:} Because Messages are containers of Fields, \emph{Assemble} and \emph{Disassemble} operators pack constituent Fields into composite Messages or split Messages into individual Fields, while \emph{Project} and \emph{Inject} select subsets of Field slots or insert constant or derived Fields, preserving Colors.
    \item \textbf{Computational Derivation and Filtering:} \emph{Derive} is the general multi-input synthesis operator: it consumes $N \ge 1$ Fields (potentially correlated from distinct Offers) and computes $M \ge 1$ output Fields with declared target Colors (for example, combining host vehicle position, attitude quaternion, and lever-arm geometry to derive sensor position, executing a state-estimation filter update).  \emph{Filter/Validate} forwards or suppresses Messages based on deterministic predicates over payload values or provenance.
    \item \textbf{Topology, Storage, and Flow Mediation:} Operators that mutate kinematics $k$ and interaction topology while leaving $(C, L, P, \rho)$ invariant.  \emph{Fan-Out}, \emph{Fan-In}, and \emph{Temporal Join} alter routing multiplicity and correlate asynchronous Messages across bounded time windows; \emph{Stores} (latest-value latches, bounded FIFO queues, or keyed maps) introduce internal memory state; \emph{Mode Bridges} reconcile mismatched push/pull transfer mechanics; and \emph{Temporal Control} operators enforce data freshness via stale-age gates and failover multiplexers.
\end{enumerate}

\begin{table}[t]
\centering
\caption{Taxonomy of Mediation Operator families and their effects on Color $\kappa = (C, L, P, \rho, k, \pi)$.}
\label{tab:legal-moves}
\footnotesize
\begin{tabularx}{\linewidth}{@{}>{\raggedright\arraybackslash}p{2.7cm} >{\raggedright\arraybackslash}p{2.8cm} >{\raggedright\arraybackslash}p{3.1cm} >{\raggedright\arraybackslash}X@{}}
\toprule
\textbf{Move Family} & \textbf{Mediation Operators} & \textbf{Color Effect} & \textbf{Kinematic \& State Effect in $Z_b = \mathrm{RSY}(\mathcal{P})$} \\
\midrule
Stateful Recoloring & Conceptual, Logical, Platform, Runtime Binding/Marshaling & Modifies $C$, $L$, $P$ ($\Delta C$, $\Delta L$, $\Delta P$) or realization $\rho$ ($\Delta\rho$, with $(C, L, P)$ invariant) & Stateless synchronous value or realization transformation; tracks precision narrowing. \\
\addlinespace[2pt]
Structural Aggregation & Assemble, Disassemble, Project, Inject & Preserves $\kappa$ of constituent Fields & Repackages Field slots between Producer and Consumer Message structures. \\
\addlinespace[2pt]
Derivation \& Filtering & Derive / Synthesize, Filter / Validate & Produces declared target $\kappa_{\text{out}}$ (extends lineage $\pi$) & Fuses $N$ inputs into $M$ derived Fields or suppresses invalid or out-of-bounds samples. \\
\addlinespace[2pt]
Flow, Storage \& Time & Fan/Join, Store, Mode Bridge, Stale Gate & Invariant in $(C, L, P, \rho)$; modifies only $k$ & Introduces internal state $S$, decouples push/pull kinematics, and enforces freshness bounds. \\
\bottomrule
\end{tabularx}
\end{table}

\begin{figure}[t]
\centering
\resizebox{\linewidth}{!}{%
\begin{tikzpicture}[
    >=Stealth,
    font=\footnotesize\sffamily,
    offerbox/.style={draw=techNavy, fill=techNavy!8, thick, rounded corners=3pt, minimum width=3.1cm, minimum height=1.15cm, align=center},
    needbox/.style={draw=boundGreen, fill=boundGreen!10, thick, rounded corners=3pt, minimum width=3.1cm, minimum height=1.15cm, align=center},
    opbox/.style={draw=compBorder, fill=compFill, thick, rounded corners=3pt, minimum width=3.1cm, minimum height=1.15cm, align=center},
    storebox/.style={draw=poiPurple, fill=poiPurple!8, thick, rounded corners=3pt, minimum width=3.1cm, minimum height=1.15cm, align=center},
    route/.style={->, thick, draw=wireInk},
    clabel/.style={font=\scriptsize\sffamily\bfseries, text=black!75}
]
    % Producer Offers (Left Boundary, x = 0)
    \node[offerbox] (o1) at (0, 1.15) {
        \textbf{Offer 1: NavMsg}\\[0pt]
        {\scriptsize $(C_{\text{ac}}, L_{\text{wgs}}, P_{\text{f64}})$}\\[-1pt]
        {\scriptsize Mode: \texttt{PUSH @ 10Hz}}
    };
    \node[offerbox] (o2) at (0, -1.15) {
        \textbf{Offer 2: AttMsg}\\[0pt]
        {\scriptsize $(C_{\text{att}}, L_{\text{body}}, P_{\text{quat}})$}\\[-1pt]
        {\scriptsize Mode: \texttt{PUSH @ 15Hz}}
    };

    % Interior Board Column 1 (x = 4.2)
    \node[opbox] (dis) at (4.2, 1.15) {
        \textbf{1. Disassemble}\\[0pt]
        {\scriptsize Project pos slot ($\Delta c = 0$)}\\[-1pt]
        {\scriptsize Out: $(C_{\text{ac}}, L_{\text{wgs}}, P_{\text{f64}})$}
    };
    \node[opbox] (join) at (4.2, -1.15) {
        \textbf{2. Temporal Join}\\[0pt]
        {\scriptsize Window $\Delta t \le 10\,\text{ms}$}\\[-1pt]
        {\scriptsize Correlates $O_1 + O_2$}
    };

    % Interior Board Column 2 (x = 8.4)
    \node[opbox, fill=fsrBlue!8, draw=fsrBlue] (der) at (8.4, -1.15) {
        \textbf{3. Derive \& Recolor}\\[0pt]
        {\scriptsize $\Delta C, \Delta L$: Lever-Arm}\\[-1pt]
        {\scriptsize Out: $(C_{\text{sens}}, L_{\text{ned}}, P_{\text{f64}})$}
    };
    \node[opbox, fill=bsrTeal!8, draw=bsrTeal] (prec) at (8.4, 1.15) {
        \textbf{4. Platform Recolor}\\[0pt]
        {\scriptsize $\Delta P$: \texttt{f64} $\to$ \texttt{f32}}\\[-1pt]
        {\scriptsize Out: $(C_{\text{sens}}, L_{\text{ned}}, P_{\text{f32}})$}
    };

    % Interior Board Column 3 (x = 12.6)
    \node[storebox] (store) at (12.6, 1.15) {
        \textbf{5. Latest Store}\\[0pt]
        {\scriptsize Retains latest sample}\\[-1pt]
        {\scriptsize Color: $(C_{\text{sens}}, L_{\text{ned}}, P_{\text{f32}})$}
    };
    \node[storebox] (gate) at (12.6, -1.15) {
        \textbf{6. Mode Bridge \& Gate}\\[0pt]
        {\scriptsize Freshness $\le 50\,\text{ms}$}\\[-1pt]
        {\scriptsize Adapts \texttt{PUSH} $\to$ \texttt{PULL}}
    };

    % Consumer Need (Right Boundary, x = 16.8)
    \node[needbox] (need) at (16.8, -1.15) {
        \textbf{Need: SensorPos}\\[0pt]
        {\scriptsize $(C_{\text{sens}}, L_{\text{ned}}, P_{\text{f32}})$}\\[-1pt]
        {\scriptsize Mode: \texttt{PULL (On-Demand)}}
    };

    % Routes
    \draw[route] (o1) -- node[above=1pt, clabel] {\texttt{PUSH}} (dis);
    \draw[route] (o2) -- node[above=1pt, clabel] {\texttt{PUSH}} (join);
    \draw[route] (dis) -- node[right=2pt, clabel] {pos field} (join);
    \draw[route] (join) -- node[above=1pt, clabel] {Joined} (der);
    \draw[route] (der) -- node[right=2pt, clabel] {$\Delta C, \Delta L$} (prec);
    \draw[route] (prec) -- node[above=1pt, clabel] {$\Delta P$} (store);
    \draw[route] (store) -- node[right=2pt, clabel] {Cached} (gate);
    \draw[route] (gate) -- node[above=1pt, clabel] {\texttt{PULL}} (need);

    % Background Board Frame
    \begin{scope}[on background layer]
        \node[draw=black!45, dashed, thick, fill=black!2, rounded corners=5pt, fit=(dis)(join)(der)(prec)(store)(gate), inner sep=8pt,
              label={[anchor=south west, font=\footnotesize\sffamily\bfseries, text=black!75]north west:Mediation Board Play $\mathcal{P}$ (Coupled Resultant $Z_b = \mathrm{RSY}(\mathcal{P})$)}] {};
    \end{scope}
\end{tikzpicture}%
}
\caption{An example Play $\mathcal{P}$ on a Mediation Board.  Two asynchronous push-mode Producer Offers are disassembled, temporally joined, transformed across conceptual and logical axes via a lever-arm derivation, narrowed in platform precision, and buffered in a latest-value store with a freshness-gated mode bridge to satisfy a pull-mode Consumer Need.  Operator boxes display the $(C, L, P)$ projection of each Field's Color $\kappa$ for brevity; $\rho$, $k$, and $\pi$ are carried along and mutate only where an operator (e.g., the Store and Mode Bridge, which alter $k$) is applied.}
\label{fig:mediation-board}
\end{figure}

Figure~\ref{fig:mediation-board} illustrates how a Play $\mathcal{P}$ composes these operator families to resolve both $(C, L, P)$ differences and the push-to-pull kinematic gap between two Producer Offers and a Consumer Need.

\subsection{Behavioral Contracts} \label{sec:meth-contracts}

Typing alone does not establish that a Coupled Resultant behaves correctly or as desired over time.  We encode behavioral requirements as Input/Output Requirements $IOR = (OLR, IR, ITR, OR, OTR, ER)$, where the operational life requirement satisfies $OLR \in \mathbb{N}^+ \cup \{\infty\}$ and induces the timescale $TSR(IOR) = \{t \in \mathbb{N} \mid t < OLR\}$.  When $OLR = \infty$, $TSR = \mathbb{N}$.  To encode $IOR$s without loss of generality, we use a Universal Live Sequence Chart (uLSC) fragment over two boundary lifelines: the external environment ($\mathsf{Env}$) and the system ($\mathsf{Sys}$).

\begin{definition}[Universal Live Sequence Chart (uLSC)] \label{def:ulsc}
A uLSC over input alphabet $IR$ and output alphabet $OR$ is a tuple $\mathcal{L} = (\tau, \mathcal{P}_{\text{in}}, \mathcal{A}_{\text{out}})$, where $\tau \subseteq \mathbb{N}$ is the discrete operational horizon (finite, or $\tau = \mathbb{N}$ when $OLR = \infty$), $\mathcal{P}_{\text{in}} \subseteq IR^{\mathbb{N}}$ is the \emph{cold Pre-Chart} predicate specifying the set of eligible environmental input trajectories that activate the contract, and $\mathcal{A}_{\text{out}} : IR^{\mathbb{N}} \times OR^{\mathbb{N}} \to \{\mathrm{true}, \mathrm{false}\}$ is the \emph{hot Main-Chart} relation over $\omega$-trajectories, specifying admissible input--output trajectory pairs (including value tolerances, delivery deadlines, and quiet-output \emph{silence} obligations).
\end{definition}

A discrete system $Z = (S, IR, OR, \delta, \rho_{\text{out}})$ initialized at $s_0 \in S$ \emph{satisfies} a uLSC $\mathcal{L}$ on timescale $T \subseteq \tau$ ($T \neq \emptyset$), written $(Z, s_0) \models_T \mathcal{L}$, when every Pre-Chart activation forces a Main-Chart fulfillment:
\begin{equation}
\begin{aligned}
    (Z, s_0) \models_T \mathcal{L}
    \iff{}&
    \forall f \in \mathcal{P}_{\text{in}},\; \forall h \in IR^{\mathbb{N}},\\
    &\quad\big(\forall t \in \tau,\; h(t) = f(t)\big) \implies \exists g \in OR^{\mathbb{N}},\; \mathcal{A}_{\text{out}}(f, g) \wedge \big(\forall t \in T,\; y_Z(s_0, h)(t) = g(t)\big).
\end{aligned}
    \label{eq:lsc-sat}
\end{equation}
We write $\mathrm{CTL}(\mathcal{L})$ for the set of all triples $(Z, s_0, T)$ such that $(Z, s_0) \models_T \mathcal{L}$.  If an input trajectory $h$ fails to match $\mathcal{P}_{\text{in}}$ on $\tau$, the cold Pre-Chart is vacated without penalty; conversely, if the Pre-Chart matches on $\tau$ but no allowed output trajectory $g$ agrees with $y_Z(s_0, h)$ on $T$, a \emph{Hot Breach} occurs, written $(Z, s_0, f) \not\models_T \mathcal{L}$ for the activating scenario $f \in \mathcal{P}_{\text{in}}$.  Given any Wymorian Input/Output Requirement $IOR$, its \emph{characteristic uLSC} $\mathcal{L}(IOR)$ is constructed by setting $\tau = TSR(IOR)$, $\mathcal{P}_{\text{in}} = ITR$, and $\mathcal{A}_{\text{out}}(f, g) \Longleftrightarrow g \in ER(f)$ (Fig.~\ref{fig:ulsc-diagram}).  Because $OLR \in \mathbb{N}^+ \cup \{\infty\}$, this construction covers both finite horizons and the infinite horizon $\tau = TSR = \mathbb{N}$, in which $\mathcal{P}_{\text{in}} \subseteq IR^{\mathbb{N}}$ and $\mathcal{A}_{\text{out}}$ range over $\omega$-trajectories.

\begin{theorem}[Correspondence Between $IOR$s and uLSCs] \label{thm:ior-lsc}
For every Input/Output Requirement $IOR = (OLR, IR, ITR, OR, OTR, ER)$ with $OLR \in \mathbb{N}^+ \cup \{\infty\}$, discrete system $Z = (S, IR, OR, \delta, \rho_{\text{out}})$, initial state $s_0 \in S$, and nonempty evaluation timescale $T \subseteq TSR(IOR)$,
\begin{equation}
    (Z, s_0) \models_T \mathcal{L}(IOR)
    \iff (Z, s_0) \models_T IOR
    \iff (Z, s_0, T) \in \mathrm{CTL}(IOR).
    \label{eq:ior-lsc-equiv}
\end{equation}
Consequently $\mathrm{CTL}(\mathcal{L}(IOR)) = \mathrm{CTL}(IOR)$.  In particular, the equivalence holds in the infinite case $OLR = \infty$, $TSR = \mathbb{N}$, $T = \mathbb{N}$.  Moreover, a Hot Breach occurs on an input scenario $f \in ITR$ (with extension $h$ agreeing with $f$ on $TSR$) if and only if $(Z, s_0, f) \not\models_T IOR$.
\end{theorem}

\begin{proof}[Proof Sketch]
By definition of $\mathcal{L}(IOR)$, $\tau = TSR(IOR)$, $\mathcal{P}_{\text{in}} = ITR$, and $\mathcal{A}_{\text{out}}(f, g) \iff g \in ER(f)$.  Substituting these directly into~\eqref{eq:lsc-sat} yields the exact trajectory quantification of $(Z, s_0) \models_T IOR$ from Section~\ref{sec:bg-t3sd}: $T \neq \emptyset$, $T \subseteq TSR(IOR)$, and for every $f \in ITR$ and extension $h$ agreeing with $f$ on $TSR(IOR)$, there exists $g \in ER(f)$ matching $y_Z(s_0, h)(t)$ at all $t \in T$.  Membership $(Z, s_0, T) \in \mathrm{CTL}(IOR)$ is synonymous with that quantification, so the three sides of~\eqref{eq:ior-lsc-equiv} coincide and $\mathrm{CTL}(\mathcal{L}(IOR)) = \mathrm{CTL}(IOR)$.  The argument is syntactic and never uses finiteness of $\tau$ or $T$, so it covers the infinite case: when $OLR = \infty$ we have $\tau = \mathbb{N}$, the agreement condition $\forall t \in \tau,\, h(t) = f(t)$ forces $h = f$, $g \in ER(f) \subseteq OR^{\mathbb{N}}$ ranges over entire $\omega$-trajectories, and the universal quantifier over $t \in T = \mathbb{N}$ is carried identically on both sides of~\eqref{eq:ior-lsc-equiv}.  Defining $\mathcal{A}_{\text{out}}$ via full envelope membership $g \in ER(f)$ is essential: a pointwise projection requiring only $y_Z(s_0, h)(t) \in \{g(t) \mid g \in ER(f)\}$ at each tick $t$ independently is strictly weaker whenever $ER(f)$ couples outputs across ticks.  Negating~\eqref{eq:lsc-sat} immediately yields the Hot Breach condition $(Z, s_0, f) \not\models_T \mathcal{L}$.
\end{proof}

\begin{figure}[t]
\centering
\begin{tikzpicture}[
    >=Stealth,
    font=\footnotesize\sffamily,
    lifelinehead/.style={draw=compBorder, fill=compFill, thick, rounded corners=3pt, minimum width=2.8cm, minimum height=0.68cm, align=center, font=\footnotesize\sffamily\bfseries},
    coldframe/.style={draw=coldBlue, dashed, thick, fill=coldBlue!4, rounded corners=4pt},
    hotframe/.style={draw=hotRed, solid, thick, fill=hotRed!4, rounded corners=4pt},
    lsccond/.style={draw=black!30, fill=white, rounded corners=3pt, font=\scriptsize\sffamily, inner sep=3.5pt, align=center}
]
    % Lifeline Heads
    \node[lifelinehead] (env) at (0, 0) {$\mathsf{Env}$};
    \node[lifelinehead] (sys) at (8.4, 0) {$\mathsf{Sys}$};

    % Pre-Chart (Cold Trigger Region): y from -0.85 to -2.75
    \node[coldframe, minimum width=10.2cm, minimum height=1.90cm, anchor=north] (pre_box) at (4.2, -0.85) {};
    \node[font=\scriptsize\sffamily\bfseries, text=coldBlue, anchor=south west] at ([yshift=1pt]pre_box.north west)
      {Pre-Chart: Cold Trigger ($f \in ITR$)};

    % Main-Chart (Hot Guarantee Region): y from -3.30 to -5.70
    \node[hotframe, minimum width=10.2cm, minimum height=2.40cm, anchor=north] (main_box) at (4.2, -3.30) {};
    \node[font=\scriptsize\sffamily\bfseries, text=hotRed, anchor=south west] at ([yshift=1pt]main_box.north west)
      {Main-Chart: Hot Guarantee ($g \in ER(f)$)};

    % Vertical Lifeline Rays
    \draw[dashed, thick, black!45] (env.south) -- (0, -5.95);
    \draw[dashed, thick, black!45] (sys.south) -- (8.4, -5.95);

    % Cold Input Message (Env -> Sys)
    \draw[-{Stealth}, dashed, thick, draw=coldBlue] (0, -1.45) --
      node[above=1.5pt, font=\scriptsize\sffamily, text=coldBlue] {\texttt{OfferMsg}$(f(t)) : (C_{\text{in}}, L_{\text{in}}, P_{\text{in}})$} (8.4, -1.45);

    % Cold Guard Condition
    \node[lsccond, draw=coldBlue, dashed, thick] at (4.2, -2.22)
      {\textbf{Cold Condition:} $f \in ITR$ (Eligible Producer Scenario)};

    % Tick Synchronization
    \draw[dashdotted, thick, black!45] (-0.5, -3.90) --
      node[above=1.5pt, font=\scriptsize\sffamily, text=black!70] {Discrete Tick Synchronization $t \in TSR$ / Deadline Window $[t, t+\Delta]$} (8.9, -3.90);

    % Hot Output Message (Sys -> Env)
    \draw[-{Stealth}, solid, thick, draw=hotRed] (8.4, -4.52) --
      node[above=1.5pt, font=\scriptsize\sffamily, text=hotRed] {\texttt{NeedMsg}$(g(t)) : (C_{\text{out}}, L_{\text{out}}, P_{\text{out}})$} (0, -4.52);

    % Hot Envelope Condition
    \node[lsccond, draw=hotRed, solid, thick] at (4.2, -5.20)
      {\textbf{Hot Guarantee:} $g \in ER(f)$ \enspace$\wedge$\enspace $\mathrm{silenceOut}$ on quiet ticks (Violation $=$ \textbf{Hot Breach})};
\end{tikzpicture}
\caption{Diagrammatic structure of the characteristic uLSC $\mathcal{L}(IOR)$.  A match on the dashed blue Cold Pre-Chart triggers the solid red Hot Main-Chart obligation; any violation of the output value envelope, timing window, or silence constraint induces a Hot Breach.}
\label{fig:ulsc-diagram}
\end{figure}

\subsection{Winning Conditions} \label{sec:meth-winning}

We evaluate a candidate Play $\mathcal{P}$ across three tiers:
\begin{enumerate}[label=\textbf{Tier \arabic*}, leftmargin=3.4em, topsep=2pt, itemsep=1pt]
    \item \textbf{(Static Stateful-Color and Structural Compatibility):} For every required Field slot in every Consumer Need $N_j \in \mathcal{N}$, there exists a directed path through $\mathcal{P}$ originating at one or more Producer Offers $O_i \in \mathcal{O}$ along which only legal Mediation Operators are applied, Message aggregations assemble with matching slot structures, and the terminal Field's Color $\kappa$ matches the required $\kappa_{N_j}$: $(C, L, P, \rho)$ agree with $(C_{N_j}, L_{N_j}, P_{N_j}, \rho_{N_j})$, and $\pi$ satisfies the provenance predicate $\pi_{N_j}$.
    \item \textbf{(Flow and Kinematic Compatibility):} Along every route in $\mathcal{P}$, the kinematic component $k$ is mutually compatible (every push output connects to a push-capable input, and every pull input connects to a pull-capable output, or a Mode Bridge mediates the transition), staleness/freshness bounds are met by operator and Route Carrier latencies, all storage capacities are strictly positive, and all temporal operators share a consistent time base.
    \item \textbf{(Behavioral Satisfaction):} Let $Z_b = \mathrm{RSY}(\mathcal{P})$ be the Coupled Resultant induced by the Play.  Given the problem's $IOR$ and its characteristic uLSC $\mathcal{L}(IOR)$ (Definition~\ref{def:ulsc}), the Coupled Resultant satisfies the behavioral contract over the declared operational horizon without a Hot Breach:
    \[
        (Z_b, s_0) \models_{TSR} \mathcal{L}(IOR).
    \]
    By Theorem~\ref{thm:ior-lsc} this is equivalent to $(Z_b, s_0) \models_{TSR} IOR$ and to $(Z_b, s_0, TSR) \in \mathrm{CTL}(IOR)$, including for $OLR = \infty$.
\end{enumerate}

A Play $\mathcal{P}$ \emph{wins} the Integration Game if and only if it satisfies Tiers~1, 2, and~3 simultaneously.  By Proposition~\ref{prop:wymore-iso} and Theorem~\ref{thm:ior-lsc}, a winning Play constitutes a well-formed Color-typed coupling recipe over $TYR_{\mathcal{K}}$ whose Coupled Resultant obeys $(Z_b, s_0) \models_{TSR} \mathcal{L}(IOR)$ over the verified trajectory suite $ITR_{\text{test}} \subseteq ITR$, and across all admissible trajectories within any symbolically certified contract fragment (Section~\ref{sec:meth-rl}).

\subsection{Cost Functional and Trade-Offs} \label{sec:meth-cost}

Winning Plays can still differ dramatically in footprint, latency, and fidelity.  We therefore define an Integration Cost Functional $\mathcal{J}(\mathcal{P})$ that both supports Wymorian $CR$/$PR$/$TR$ trade offs and supplies guidance for RL: frontier potential $\mathcal{J}_{\text{tier1,2}}$ guides assembly of sparse valid graphs (Tiers~1--2), $C_{\text{breach}}$ is a temporal/value guide over traces (Tier~3), $C_{\text{static}}$ versus $C_{\text{dyn}}$ mirrors design-time Cost Requirements versus runtime Performance Requirements, and the piecewise fitness $\mathcal{S}$ preserves exact cost differences among feasible Plays while compressing $\Omega$-scale penalties without saturating gradients.

\begin{enumerate}[topsep=2pt, itemsep=1pt, leftmargin=1.8em]
    \item \textbf{Static Construction Cost $C_{\text{static}}(\mathcal{P})$ ($CR$):} Quantifies the structural and resource footprint of the Mediation Board:
    \begin{equation}
        C_{\text{static}}(\mathcal{P}) = \sum_{v \in \mathcal{V}_{\text{op}}} c_{\text{op}}(v) + \sum_{s \in \mathcal{V}_{\text{store}}} c_{\text{mem}}(\mathrm{cap}_s, \mathrm{ret}_s) + \sum_{r \in \mathcal{E}_{\text{route}}} c_{\text{carrier}}(r),
        \label{eq:c-static}
    \end{equation}
    where $\mathcal{V}_{\text{store}} \subseteq \mathcal{V}_{\text{op}}$ are the Store operators.  In general, $(c_{\text{op}}, c_{\text{mem}}, c_{\text{carrier}})$ are normalized non-negative resource weights derived from target-platform profiling or catalog specifications (e.g., WCET/FLOPs per invocation for $c_{\text{op}}$, allocated buffer bytes $\times$ retention ticks for $c_{\text{mem}}$, and bus bandwidth or IPC marshaling overhead for $c_{\text{carrier}}$).  In our implementation, operators are assigned $c_{\text{op}}$ by family: stateless recoloring/projection in $[0.1, 0.3]$, stateful stores/bridges in $[0.3, 0.5]$ plus $c_{\text{mem}}$ proportional to capacity, and multi-input derivations or heavy batch solvers in $[0.5, 1.5]$.

    \item \textbf{Dynamic Operational Cost $C_{\text{dyn}}(\mathcal{P}, f)$ ($PR$):} Quantifies trajectory-level quality-of-service degradation over an input scenario $f \in ITR$ across time $t \in TSR$ (with temporal discount factor $\gamma \in (0, 1]$):
    \begin{equation}
        C_{\text{dyn}}(\mathcal{P}, f) = \sum_{t \in TSR} \gamma^t \Big[ w_{\text{lat}}\,\mathrm{Latency}(t) + w_{\text{loss}}\,\mathrm{InfoLoss}(t) + w_{\text{drop}}\,\mathbf{1}_{\{\mathrm{Code}(t) \neq \mathsf{Success}\}} \Big],
        \label{eq:c-dyn}
    \end{equation}
    capturing end-to-end transport latency, precision loss from narrowing platform conversions, and runtime flow disruptions (e.g., stale sample drops).

    \item \textbf{Behavioral Breach Penalty $C_{\text{breach}}(\mathcal{P}, f)$:} Quantifies departure from the required behavioral envelope $ER(f)$:
    \begin{equation}
        C_{\text{breach}}(\mathcal{P}, f) = \Omega \cdot \mathbf{1}_{\mathsf{HotBreach}}(\mathcal{P}, f) + \sum_{t \in TSR} \inf_{g \in ER(f)} d_{\text{val}}\big(y_{Z_b}(t),\, g(t)\big),
        \label{eq:c-breach}
    \end{equation}
    where $\Omega > 0$ is a categorical breach margin triggered upon any uLSC Hot Breach, and $d_{\text{val}}$ is a continuous distance metric measuring how far the generated output $y_{Z_b}(t)$ deviates from the admissible envelope $ER(f)$.
\end{enumerate}

Combining these terms with the structural and kinematic frontier potentials over a distribution of operational scenarios $\mathcal{D}(ITR)$ yields the composite Integration Cost Functional:
\begin{equation}
    \mathcal{J}(\mathcal{P}) = \mathcal{J}_{\text{tier1,2}}(\mathcal{P}) + C_{\text{static}}(\mathcal{P}) + \mathbb{E}_{f \sim \mathcal{D}(ITR)}\!\Big[ C_{\text{dyn}}(\mathcal{P}, f) + C_{\text{breach}}(\mathcal{P}, f) \Big],
    \label{eq:j-total}
\end{equation}
where $\mathcal{J}_{\text{tier1,2}} = \mathcal{J}_{\text{tier1}} + \mathcal{J}_{\text{tier2}}$ is the open-frontier potential, defined explicitly as
\begin{equation}
    \mathcal{J}_{\text{tier1,2}}(\mathcal{P}) = \sum_{N_j \in \mathcal{N}_{\text{open}}(\mathcal{P})} \Big( w_{\text{unconnected}} + \min_{O \in \mathcal{F}_{\text{live}}(\mathcal{P})} \Delta_{\kappa}(O, N_j) \Big).
    \label{eq:j-frontier}
\end{equation}
Here $\mathcal{N}_{\text{open}}(\mathcal{P}) \subseteq \mathcal{N}$ is the set of Needs not yet satisfied at Tiers~1 and~2, $\mathcal{F}_{\text{live}}(\mathcal{P})$ is the \emph{live offer frontier} (i.e., the set of Offers reachable via directed paths from boundary Producers) and $w_{\text{unconnected}}$ is an open-port penalty.  The term $\Delta_{\kappa}(O, N_j)$ is the \emph{Color distance}: the sum $\Delta_{\text{tier1}} + \Delta_{\text{tier2}}$ of a Tier~1 distance over the axes $(C, L, P, \rho, \pi)$ and a Tier~2 kinematic mismatch over $k$.  By convention the minimum over an empty frontier equals a fixed maximal distance $\Delta_{\max}$.  Tier~3 contributes to $\mathcal{J}$ only through $C_{\text{breach}}$ and is \emph{not} a frontier term.

\paragraph{Fitness.} The fitness $\mathcal{S}(\mathcal{P})$ is a strictly decreasing, $C^1$-smooth map of a soft cost $\mathrm{softJ}(\mathcal{P}) = \mathcal{J}(\mathcal{P}) / J_{\text{ref}}$ (the cost $\mathcal{J}$ normalized by a fixed reference scale $J_{\text{ref}} > 0$) into $(0, S_{\max}]$, with a piecewise form that is linear up to a knee $x_k \in (0, S_{\max})$ and rational beyond it:
\begin{equation}
    \mathcal{S}(\mathcal{P}) = \sigma\big(\mathrm{softJ}(\mathcal{P})\big), \qquad
    \sigma(x) =
    \begin{cases}
        S_{\max} - x, & 0 \le x \le x_k, \\[3pt]
        \dfrac{(S_{\max} - x_k)^2}{x + S_{\max} - 2x_k}, & x > x_k.
    \end{cases}
    \label{eq:fitness-S}
\end{equation}
The two branches agree in value ($S_{\max} - x_k$) and slope ($-1$) at $x = x_k$, so $\sigma$ is $C^1$, strictly decreasing, and bounded below by $0$ without ever reaching it; fitness therefore preserves a nonvanishing gradient across arbitrarily poor candidates.

To guarantee that $\mathcal{J}(\mathcal{P})$ induces a well-posed value ordering over candidate architectures, we desire that it aligns with some expectations we have about such a cost functional should behave.  To capture these intuitive expectations, we define seven axioms:
\begin{enumerate}[label=\textbf{(A\arabic*)}, leftmargin=2.4em, topsep=2pt, itemsep=1pt]
    \item \textbf{Feasibility Dominance:} Every winning Play ($\mathcal{J}_{\text{tier1,2}}=C_{\text{breach}}=0$) achieves strictly lower cost $\mathcal{J}$ and strictly higher fitness $\mathcal{S}$ than any losing Play.
    \item \textbf{Only Count Connected Moves:} Placing an interior Mediation Operator disconnected from Producers leaves $\mathcal{F}_{\text{live}}(\mathcal{P})$ unchanged and strictly increases $\mathcal{J}$ by $c_{\text{op}}(v) > 0$.
    \item \textbf{Prefer Useful Moves:} Resolving each mismatch axis of $\kappa$ ($\Delta C, \Delta L, \Delta P, \Delta\rho, \Delta\pi$, or the kinematic transfer mode in $k$) along $\mathcal{F}_{\text{live}}$ strictly decreases $\mathcal{J}_{\text{tier1,2}}+C_{\text{static}}$ and increases $\mathcal{S}$.
    \item \textbf{Connection Continuity:} Connecting an Offer $O^* \in \mathcal{F}_{\text{live}}$ that realizes the frontier minimum, with $\Delta_{\kappa}(O^*, N_j) = 0$, to an open Need $N_j$ reduces $\mathcal{J}_{\text{tier1,2}}$ by exactly $w_{\text{unconnected}}$ without introducing a discontinuity.
    \item \textbf{Rank Breaches by Severity:} $C_{\text{breach}}$ is uncapped, and both $C_{\text{breach}}$ and $-\mathcal{S}$ are strictly monotonic with respect to chart-cut progress deficiency, deadline overrun, value error $d_{\text{val}}$, and silence violations.
    \item \textbf{Prefer Better QoS:} Among winning Plays of equal static cost, increasing operator latency, precision loss, or flow disruption events strictly increases $C_{\text{dyn}}$ and $\mathcal{J}$, and strictly decreases $\mathcal{S}$.
    \item \textbf{Prefer Simpler Boards:} Appending a redundant identity operator or a higher-cost Route Carrier to a winning Play strictly increases $C_{\text{static}}$ and $\mathcal{J}$, and strictly decreases $\mathcal{S}$.
\end{enumerate}
Axiom~(A1) holds for boards with at most $V_{\max}$ operators and $E_{\max}$ routes whenever $\Omega > \overline{C}_{\text{static}} + \overline{C}_{\text{dyn}}$, where the bars denote worst-case static and discounted dynamic cost over that bound (or $(1-\gamma)^{-1}$ for infinite horizons).  In our experiments feasible Plays satisfy $\mathcal{J}_{\text{win}} \le 6.30$ with bounded-graph cost $< 100$, so $\Omega = 10^{4}$ is any margin above that bound and yields the observed separation $\max\mathcal{J}_{\text{win}} < \min\mathcal{J}_{\text{lose}}$.

Other implementations may desire a functional that behaves according to different axioms.  We present our choices of axioms here to clarify the methodology and our implementation, not to suggest that these are the correct or only choices.

\subsection{Reinforcement Learning and Decidability} \label{sec:meth-rl}

Because the Integration Game provides a closed-form state space, action constraints, and a well-defined objective $\mathcal{J}(\mathcal{P})$, it supports two complementary decision-making regimes (Fig.~\ref{fig:rl-architecture}).

\paragraph{Offline Mode (Architecture Synthesis).}
In offline synthesis, the learning agent acts as an automated systems architect that constructs a winning Play $\mathcal{P}$ prior to deployment.  We formulate board construction as an action-masked Markov Decision Process $(S, A, P, R, \gamma)$ whose state observations summarize open Consumer Needs, the frontier potentials $\mathcal{J}_{\text{tier1}}$ and $\mathcal{J}_{\text{tier2}}$, static cost $C_{\text{static}}$, fitness $\mathcal{S}$, and the behavioral monitor state $u_t$.  Tier~3 enters the observation only through the monitor state $u_t$ and the fitness $\mathcal{S}$, never as a frontier potential.  Legal actions select Mediation Operators from $TYR_{\mathcal{K}}$, wire compatible Offer--Need port pairs, or finalize the architecture; invalid actions (e.g., connecting ports with mismatched Colors without an intervening Mediation Operator) are masked structurally~\cite{Huang2022Masking}.  To prevent sparse-reward stagnation during multi-hop graph assembly, each construction step receives a potential-based shaping reward~\cite{NgHaradaRussell1999}:
\begin{equation}
    r_{\text{shape}}(s_k, a_k, s_{k+1}) = -c_{\text{op}}(a_k) + \big(\Phi_{\kappa}(s_k) - \Phi_{\kappa}(s_{k+1})\big),
    \label{eq:pbrs}
\end{equation}
where the potential $\Phi_{\kappa}(s) = \mathcal{J}_{\text{tier1}} + \mathcal{J}_{\text{tier2}} = \mathcal{J}_{\text{tier1,2}}$ is the open-frontier potential~\eqref{eq:j-frontier} of the partial Play at state $s$, measuring the aggregate Color distance $\Delta_\kappa$ between the live offer frontier $\mathcal{F}_{\text{live}}$ and the remaining unsatisfied Needs.  Once a Play satisfies Tiers~1 and~2, trajectory rollouts over $f \sim \mathcal{D}(ITR)$ yield the terminal return $R_{\text{term}} = -\mathcal{J}(\mathcal{P})$.  To support zero-shot transfer across scenarios with distinct operator catalogs, policy networks can condition action logits on relational port-compatibility features $\phi(s, a)$ that encode which discrepancy in $\kappa$ a candidate move resolves.

\paragraph{Online Mode (Adaptive Runtime Mediation).}
Once a Mediation Board topology $\mathcal{P}$ is deployed, an online controller can regulate dynamic flow parameters across ticks $t \in TSR$ in response to non-stationary arrival rates, bursty emissions, or degraded communication links.  The runtime state $s_t \in \mathcal{S}_{\text{dyn}}$ tracks buffer occupancies, cached sample ages, link health indicators, and recent transfer codes; actions $a_t \in \mathcal{A}_{\text{dyn}}$ modulate polling intervals, queue eviction rules, or active channel selection across redundant failover paths to maximize the step reward $r_t = -\big(c_{\text{dyn}}(s_t, a_t) + c_{\text{breach}}(s_t, a_t)\big)$.

\begin{figure}[t]
\centering
\begin{tikzpicture}[
    >=Stealth,
    font=\footnotesize\sffamily,
    agentbox/.style={draw=accentAmber!90!black, fill=accentAmber!10, very thick, rounded corners=4pt, minimum width=4.5cm, minimum height=1.20cm, align=center},
    envbox/.style={draw=techNavy, fill=techNavy!8, thick, rounded corners=4pt, minimum width=5.1cm, minimum height=1.15cm, align=center},
    rmbox/.style={draw=hotRed!85!black, fill=hotRed!6, thick, rounded corners=4pt, minimum width=5.1cm, minimum height=1.15cm, align=center},
    arrow/.style={->, thick, draw=black!75}
]
    % Left: RL Policy
    \node[agentbox] (agent) at (0, 0) {
        \textbf{RL Agent Policy $\pi_\theta(\tilde{s})$}\\[0pt]
        {\scriptsize Offline: Graph Synthesis (Masked PPO)}\\[-1pt]
        {\scriptsize Online: Flow Control (Actor--Critic)}
    };

    % Top Right: Mediation Board Environment
    \node[envbox] (board) at (8.2, 1.10) {
        \textbf{Mediation Board $\mathcal{B}$}\\[0pt]
        {\scriptsize Coupled Resultant $Z_b = \mathrm{RSY}(\mathcal{P})$}\\[-1pt]
        {\scriptsize State $s \in \mathcal{S}$, Color Potential $\Phi_{\kappa}(s)$}
    };

    % Bottom Right: uLSC Reward Machine
    \node[rmbox] (rm) at (8.2, -1.10) {
        \textbf{uLSC Reward Machine $\mathcal{M}_{\text{LSC}}$}\\[0pt]
        {\scriptsize State $u \in \{\mathsf{Idle}, \mathsf{Pre}, \mathsf{Main}, \mathsf{Done}, \mathsf{Breach}\}$}\\[-1pt]
        {\scriptsize Monitors $f \in ITR \implies y_{Z_b} \in ER(f)$}
    };

    % Connections
    \draw[arrow] (agent.north) |- node[pos=0.75, above=1.5pt, font=\scriptsize] {Action $a$ (Place Operator / Wire / Tune Flow)} (board.west);
    \draw[arrow] (board.south) -- node[right=3pt, font=\scriptsize] {Boundary Trace $(f(t), y_{Z_b}(t))$} (rm.north);
    \draw[arrow] (rm.west) -| node[pos=0.25, below=1.5pt, font=\scriptsize] {Augmented State $\tilde{s} = (s, u)$ \& Reward $r = -\mathcal{J}$} (agent.south);
    \draw[arrow, dashed, draw=techNavy!80!black] ([yshift=-0.22cm]board.west) -- ++(-1.1,0) |- node[pos=0.75, above=1.5pt, font=\scriptsize, text=techNavy!90!black] {Shaping $\Delta\Phi_{\kappa}$} (agent.east);
\end{tikzpicture}
\caption{Closed-loop Reinforcement Learning architecture for the proposed \emph{Integration Game}.  Compiling the uLSC $\mathcal{L}(IOR)$ into a finite-state Reward Machine $\mathcal{M}_{\text{LSC}}$ exposes the automaton state $u$ to the agent, restoring the Markov property for trajectory-level behavioral contracts; the shaping term $\Delta\Phi_{\kappa}$ is the change in the open-frontier potential of the partial Play.}
\label{fig:rl-architecture}
\end{figure}

\paragraph{Reward Machines.}
A fundamental challenge in applying RL to Tier~3 behavioral satisfiability is that trajectory obligations $y_{Z_b} \in ER(f)$ are inherently \emph{non-Markovian}: whether an output at tick $t$ is admissible depends on the history of environmental triggers and prior system emissions over $[0, t]$.  To restore the Markov property, we compile the uLSC $\mathcal{L}(IOR)$ into a finite-state \emph{Reward Machine}~\cite{ToroIcarte2022}:
\begin{equation}
    \mathcal{M}_{\text{LSC}} = (U, u_0, \Sigma, \delta_u, \sigma_r),
    \label{eq:reward-machine}
\end{equation}
where $U = \{\mathsf{Idle}, \mathsf{PreMatching}, \mathsf{MainActive}, \mathsf{SilenceCheck}, \mathsf{Completed}, \mathsf{Breach}\}$ is the set of chart monitor states of the obligation automaton, $\Sigma$ is the alphabet of observable boundary events (Offer arrivals, Need emissions, and quiet ticks), $\delta_u : U \times \Sigma \to U$ advances the chart state as the trajectory unfolds, and $\sigma_r : U \times \Sigma \to \mathbb{R}$ emits dense progress rewards as Pre-Chart triggers and Main-Chart milestones are fulfilled (or emits a breach penalty upon an absorbing transition to $\mathsf{Breach}$).  Augmenting the environment state with the monitor state, $\tilde{s}_t = (s_t, u_t) \in \mathcal{S} \times U$, restores the Markov property over the operational horizon.  When $OLR = \infty$ the monitor runs indefinitely and $\mathsf{Completed}$ is reached only per activation; the absorbing $\mathsf{Breach}$ state remains well defined because a Hot Breach is witnessed at a finite prefix.

\paragraph{Decidability.}
Whether a Play satisfies its behavioral contract is decidable under finiteness assumptions that realistic MOSA technologies typically meet:

\begin{proposition}[Finite Technologies] \label{prop:finite-decidable}
Suppose every component and Mediation Operator in $TYR_{\mathcal{K}}$ has a finite-state realization over finite platform alphabets $A_{\kappa}$, and suppose the contract $\mathcal{L}(IOR)$ is recognized by a finite-state monitor $\mathcal{M}_{\text{LSC}}$ (for $OLR < \infty$) or by an $\omega$-automaton monitor (for $OLR = \infty$).  Then for every Play $\mathcal{P}$, the Coupled Resultant $Z_b = \mathrm{RSY}(\mathcal{P})$ is a finite-state Moore machine, and Hot Breach reachability in the product $Z_b \otimes \mathcal{M}_{\text{LSC}}$ is decidable, even for $OLR = \infty$.
\end{proposition}

\begin{proof}[Proof Sketch]
Synchronous composition of finitely many finite-state Moore components over finite alphabets is a finite-state Moore machine whose state set is at most the product of the component state sets, so $Z_b$ is finite-state.  The product of $Z_b$ with the finite-state (or $\omega$-automaton) monitor $\mathcal{M}_{\text{LSC}}$ is again finite.  Reachability of the absorbing $\mathsf{Breach}$ state is graph reachability in this finite product, and for $\omega$-automaton monitors, acceptance of a breaching run reduces to emptiness checking, which is decidable.  For $OLR = \infty$ the horizon $T = \mathbb{N}$ is unbounded, but a Hot Breach of a value, deadline, or silence obligation is witnessed by a finite prefix, so reachability of $\mathsf{Breach}$ in the finite product settles the question without enumerating infinite trajectories.  By Theorem~\ref{thm:ior-lsc}, Breach-freedom of the product is equivalent to $(Z_b, s_0) \models IOR$.
\end{proof}

The product construction is exponential in the number of components, and the alphabets $A_{\kappa}$ are combinatorially large: a single \texttt{binary64} Field already admits $2^{64}$ values, and Messages multiply these.  Explicit enumeration is therefore infeasible.  Evaluating a candidate on a finite test suite $ITR_{\text{test}} \subseteq ITR$ is cheap but carries the classic verification hazard of \emph{coincidence-point overfitting}: a flawed Mediation Operator may agree with the true specification on sampled inputs while violating it elsewhere (for example, a unit conversion that leaves numerical values unchanged happens to be exact at $-40^\circ\text{F} = -40^\circ\text{C}$, and a non-commutative kinetic energy operator $\frac{1}{2} m v^2$ with swapped inputs is exact whenever $m = v$).  We therefore pair concrete trajectory evaluation with symbolic verification in a Counterexample-Guided Inductive Synthesis (CEGIS) loop~\cite{SolarLezama2008Sketch}, an instance of the broader Syntax-Guided Synthesis (SyGuS) paradigm~\cite{AlurEtAl2013SyGuS}.  To keep the symbolic queries decidable, elementary data transformations are abstracted over idealized reals $\mathbb{R}$ and restricted to quantifier-free fragments: linear real arithmetic ($\mathsf{QF\_LRA}$), Tarski-decidable polynomial real arithmetic ($\mathsf{QF\_NRA}$), and mixed linear integer/real arithmetic for bounded model checking over discrete schedules ($\mathsf{QF\_LIRA}$).  Whenever a synthesized Play passes $ITR_{\text{test}}$, a symbolic query asks whether the abstract contract admits a counterexample input trajectory $f_{\text{cex}}$.  An unsatisfiable verdict certifies correctness over the modeled fragment, whereas a satisfiable assignment yields a concrete counterexample that is appended to $ITR_{\text{test}}$ to refute the flawed architecture and guide subsequent synthesis.
